% BEGIN GENERATED ARTIFACT: def:elementary-disjunction-propositional-logic
\begin{definitionbox}{Definition (Elementary Disjunction)}
\begin{definition}[Elementary Disjunction]
\label{def:elementary-disjunction-propositional-logic}
An \emph{elementary disjunction} is a finite disjunction of
\hyperref[def:literal-propositional-logic]{literals}. That is, it is a
\hyperref[def:clause-propositional-logic]{clause} of the form
\[
L_1\lor L_2\lor\cdots\lor L_k,
\]
where each \(L_i\) is a literal. The one-literal case is permitted.
\end{definition}
\end{definitionbox}

\begin{remark*}[Standard quantified statement]
\[
L_1\lor L_2\lor\cdots\lor L_k,
\qquad
L_i\text{ is a literal for each }i.
\]
\end{remark*}

\begin{remark*}[Interpretation]
Elementary disjunctions are the clause-level building blocks used in
conjunctive normal form. They record finite literal alternatives.
\end{remark*}

\begin{example*}[Examples]
\[
P,\qquad P\lor\neg Q,\qquad \neg P\lor R\lor S
\]
are elementary disjunctions.
\end{example*}

\begin{dependencies}
\begin{itemize}
  \item \hyperref[def:clause-propositional-logic]{Clause}
  \item \hyperref[def:literal-propositional-logic]{Literal}
\end{itemize}
\end{dependencies}
% END GENERATED ARTIFACT: def:elementary-disjunction-propositional-logic

% BEGIN GENERATED ARTIFACT: def:maxterm-propositional-logic
\begin{definitionbox}{Definition (Maxterm)}
\begin{definition}[Maxterm]
\label{def:maxterm-propositional-logic}
Let \(P_1,\ldots,P_n\) be a finite list of propositional variables. A
\emph{maxterm} relative to \(P_1,\ldots,P_n\) is an
\hyperref[def:elementary-disjunction-propositional-logic]{elementary disjunction}
that contains exactly one of \(P_i\) or \(\neg P_i\) for each
\(i=1,\ldots,n\).
\end{definition}
\end{definitionbox}

\begin{remark*}[Standard quantified statement]
\[
M=L_1\lor\cdots\lor L_n,
\]
where for each \(i\), exactly one of \(P_i\) or \(\neg P_i\) appears among the
literals \(L_1,\ldots,L_n\).
\end{remark*}

\begin{remark*}[Interpretation]
A maxterm relative to \(P_1,\ldots,P_n\) specifies one literal-level
alternative for every variable in the list. It is false under exactly the
assignment that makes each listed literal false.
\end{remark*}

\begin{example*}[Examples]
Relative to \(P,Q,R\), the formula
\[
P\lor\neg Q\lor R
\]
is a maxterm. Relative to \(P,Q,R\), the formula \(P\lor R\) is not a maxterm
because it omits both \(Q\) and \(\neg Q\).
\end{example*}

\begin{dependencies}
\begin{itemize}
  \item \hyperref[def:elementary-disjunction-propositional-logic]{Elementary Disjunction}
  \item \hyperref[def:formula-variable-set-propositional-logic]{Variable Set of a Formula}
\end{itemize}
\end{dependencies}
% END GENERATED ARTIFACT: def:maxterm-propositional-logic

% BEGIN GENERATED ARTIFACT: def:conjunctive-normal-form-propositional-logic
\begin{definitionbox}{Definition (Conjunctive Normal Form)}
\begin{definition}[Conjunctive Normal Form]
\label{def:conjunctive-normal-form-propositional-logic}
A formula is in \emph{conjunctive normal form}, abbreviated \emph{CNF}, when it
is a finite conjunction of
\hyperref[def:clause-propositional-logic]{clauses}. Equivalently, it has the
shape
\[
C_1\land C_2\land\cdots\land C_m,
\]
where each \(C_j\) is a clause.
\end{definition}
\end{definitionbox}

\begin{remark*}[Standard quantified statement]
\[
\varphi\text{ is in CNF}
\Longleftrightarrow
\varphi=C_1\land C_2\land\cdots\land C_m
\]
for finitely many clauses \(C_1,\ldots,C_m\).
\end{remark*}

\begin{remark*}[Interpretation]
CNF is a syntactic form used to expose clause-by-clause constraints. Replacing
a formula by a
\hyperref[def:logical-equivalence-propositional-logic]{logically equivalent}
CNF formula turns global validity and
\hyperref[def:contradiction-propositional-logic]{contradiction} questions into
questions about the interaction of its clauses.
\end{remark*}

\begin{example*}[Examples]
\[
(P\lor Q)\land(\neg P\lor R)\land S
\]
is in CNF because each conjunct is a clause.
\end{example*}

\begin{dependencies}
\begin{itemize}
  \item \hyperref[def:clause-propositional-logic]{Clause}
  \item \hyperref[def:logical-equivalence-propositional-logic]{Logical Equivalence}
  \item \hyperref[def:contradiction-propositional-logic]{Contradiction}
\end{itemize}
\end{dependencies}
% END GENERATED ARTIFACT: def:conjunctive-normal-form-propositional-logic

% BEGIN GENERATED ARTIFACT: thm:existence-conjunctive-normal-form-propositional-logic
\begin{theorembox}{Theorem (Existence of Conjunctive Normal Form)}
\begin{theorem}[Existence of Conjunctive Normal Form]
\label{thm:existence-conjunctive-normal-form-propositional-logic}
For every formula \(\varphi\in\WFF\), there exists a formula
\(\psi\in\WFF\) such that \(\psi\) is in
\hyperref[def:conjunctive-normal-form-propositional-logic]{conjunctive normal form}
and
\[
\varphi\equiv\psi.
\]
\noindent\hyperref[prf:existence-conjunctive-normal-form-propositional-logic]{Go to proof.}
\end{theorem}
\end{theorembox}
\begin{remark*}[Standard quantified statement]
\[
(\forall\varphi\in\WFF)(\exists\psi\in\WFF)
\bigl(\psi\text{ is in CNF}\land \varphi\equiv\psi\bigr).
\]
\end{remark*}

\begin{remark*}[Predicate reading]
The predicate form records the relation or property named by the statement in
the displayed logical context.
\end{remark*}


\begin{remark*}[Interpretation]
One first converts \(\varphi\) to negation normal form, then repeatedly applies
distributive laws to move disjunctions below conjunctions. Replacement of
logical equivalents preserves equivalence through every local rewrite.
\end{remark*}

\begin{dependencies}
\begin{itemize}
  \item \hyperref[def:conjunctive-normal-form-propositional-logic]{Conjunctive Normal Form}
  \item \hyperref[def:negation-normal-form-propositional-logic]{Negation Normal Form}
  \item \hyperref[thm:distributive-laws-propositional-logic]{Distributive Laws for Propositional Logic}
  \item \hyperref[thm:replacement-of-logical-equivalents-propositional-logic]{Replacement of Logical Equivalents}
  \item \hyperref[thm:existence-negation-normal-form-propositional-logic]{Existence of Negation Normal Form}
\end{itemize}
\end{dependencies}
% END GENERATED ARTIFACT: thm:existence-conjunctive-normal-form-propositional-logic
